\subsection{Contrapositive Proof}
According to the the previous chapter, one may knows that $P \Rightarrow Q \equiv \neg Q \Rightarrow \neg P$ where $P$ and $Q$ are two statements. 
These two statements share the same true table. To prove that $P \Rightarrow Q$ is true, it suffices to instead prove that $\neg Q \Rightarrow \neg P$ is true. 

\subsubsection{Outline for Contrapositive Proof}
\begin{proposition}
    If P, then Q.
    \begin{proof}
        Suppose $\neg Q$
        $$\cdots$$
        $$\cdots$$
        Therefore $\neg P$
    \end{proof}
\end{proposition}

\begin{example}
    Prove the following proposition.
    \begin{proposition}
        Suppose $x, y \in \Z$. If $5 \nmid xy$, then $5 \nmid x$ and $5 \nmid y$.
    \end{proposition}
    Suppose $\neg \left(5 \nmid x \wedge 5 \nmid y\right)$, it equals to $(5 \mid x) \vee (5 \mid y)$. 
    
    We consider these possibilities separately.
    
    \textbf{Case 1}:$5 \mid x$
    \begin{align*}
        &x = 5a, a \in \Z \\
        &\Rightarrow xy = (5a)y = 5(ay)
    \end{align*}
    Therefore $5 \mid xy$ by the definition.
    
    \textbf{Case 2}: $5 \mid y$
    \begin{align*}
        &y = 5a, a \in \Z \\
        &\Rightarrow xy = x(5a) = 5(xa)
    \end{align*}
    Therefore $5 \mid xy$ by the definition.

    The above cases show that $5 \mid xy$, so it is not true that $5 \nmid xy$.
\end{example}

\subsection{Congruence of Integers}
\begin{definition}
    [Congruent Modulo]
    Given integers a and b and $n \in \N$, we say that a and b are congruent modulo n if $n \mid (a - b)$. We express this as $a \equiv b \pmod{n}$. 
    If a and b are not congruent modulo n, we write this as $a \not\equiv b \pmod{n}$.
\end{definition}

\begin{proposition}
    \[
        \forall a, b \in \Z, c \in \N, a \equiv b \pmod{c} \Rightarrow a^2 \equiv b^2 \pmod{c}
    \]
\end{proposition}
\begin{proof}
    \begin{align*}
        &a - b = ck, k \in Z \\
        &\Rightarrow (a + b)(a - b) = ck(a + b) \\ 
        &\Rightarrow a^2 - b^2 = ck(a + b) \\
        &\Rightarrow a^2 - b^2 = cd, d = k(a + b) \in N \\
        &\Rightarrow a^2 \equiv b^2 \pmod{n}
    \end{align*}
\end{proof}

\subsection{Mathematical Writing}
\begin{enumerate}
    \item Begin each sentence with a word, not a mathematical symbol. 
    \item End each sentence with a period. 
    \item Separate mathematical symbols and expressions with words. 
    \item Avoid misuse of symbols. 
    \item Avoid using unnecessary symbols. 
    \item Use the first person plural. 
    \item Use the active voice. 
    \item Explain every new symbol. 
    \item Watch out for "it". 
    \item Since, because, as, for so. 
    \item Thus, hence, therefore, consequently. 
    \item Clarity is the gold standard of mathematical writing. 
\end{enumerate}

\newpage